\section{Experiments}
\label{sec:experiments}

Lifting the closed program assumption to compile programs with free variables
introduces a potential complication. There is no guarantee a neural network will use the compiled
$\rd{k}$-linear map for its behavior. The goal of these experiments is to establish
that transformers can learn to use these compiled $\rd{k}$-linear maps. Our primary
case study trains transformers to implement the conditional reverse algorithm from
Fig. \ref{fig:conditional_reverse}, using one-layer transformers structured like Fig. \ref{fig:transformer}.
Appendix \ref{app:experiments}
contains additional details and experiments, replicating our findings across 
three other tasks. \textbf{In summary, these findings show that transformers
can learn to use the compiled $\rd{k}$-linear maps.}

% \vspace{-.5em}
\paragraph{Models}

We train four one-layer causal transformers \citep{radford2018improving} on sequences of the form 
\t{input|output}, where \t{|} is a delimiter separating the input from the correct output. 
The $\t{input}$ is a 15-character string sampled uniformly from lowercase letters in the
alphabet. If the first character in $\t{input}$ is a vowel, \t{output} is
the reversed string; otherwise, \t{output} is equal to \t{input}. Training uses teacher-forced 
next-token prediction with cross-entropy loss applied to tokens after the delimiter. 
The four transformers share the same base architecture and training procedure,
differing only in whether an additional compiled attention head is present and,
if so, which program is compiled into it.

\setlist[itemize]{topsep=4pt, itemsep=2pt, parsep=2pt, leftmargin=1.5em, rightmargin=1.5em}
\begin{itemize}[label=$\bullet$]
  \item \t{Model I} is our baseline model. It is a one-layer causal transformer with four attention
  heads, using learned token and position embeddings---essentially Fig. \ref{fig:conditional_reverse}
  without the compiled attention head. The \t{I} stands for indirect, as it is not \textit{directly} 
  programmed with $\cj$\!.
  \item \t{Model F} is our primary model. It is \t{Model I} but with an additional
  attention head compiled from an encoding of the conditional reverse algorithm
  in Fig. \ref{fig:conditional_reverse}. The \t{F} stands for full.
  \item \t{Model P} is an ablation model. It is \t{Model F} but with the additional 
  attention head compiled from an encoding of a reverse algorithm. It is an ablation
  in the sense that it compiles a simpler algorithm. The \t{P} stands for
  partial.
  \item \t{Model T} is an ablation model. It is \t{Model F} but where the compiled attention
  head is randomized. It is an ablation in the sense that it keeps the signature of the compiled
  attention head, but randomizes the compiled $\rd{k}$-linear map. The \t{T} stands for type-preserving.
\end{itemize}

\begin{wrapfigure}{r}{0.4\linewidth}
  \vspace{-3em}
  \centering
  \includegraphics[width=\linewidth]{../experiments/cond_reverse/data/plots/behavioral_dynamics_lr0.01_bs128_maxstep600.pdf}
  \caption{Behavioral dynamics.}
  \label{fig:behavioral-dynamics}
  \vspace{-3.5em}
\end{wrapfigure}


\subsection{Does compiling change learning?}

% Summary
Each model eventually learns to implement conditional reverse, but there are
differences in how.
% Explain Figure
Fig. \ref{fig:behavioral-dynamics} summarizes their early training dynamics. 
The heatmaps measure across 20 seeds the behavior of each model on the input:
\t{onmlkjihgfedcba}. The y-axis tracks across seeds whether a model outputs \t{a} in its first position, \t{b} in its second position,
and so forth---the correct output after reversal.

% What is interesting
Because input strings are drawn uniformly from lowercase letters in the alphabet, the
majority of the data do not contain a vowel in the first position. Therefore, a decent
strategy for minimizing loss is to just learn identity, and not worry about reversal. This
is essentially what \t{Model I} does. Its initial dip in loss comes from learning
identity, and then the rest of learning is spent conditionally reversing for successively
more vowels in the first position. \t{Model P} is similar. However, \t{Model F} never
transitions into an overuse of identity, and instead outright learns the task the fastest
across models, suggesting that the compiled attention head may be driving
its behavior. \t{Model T} is similar.

\subsection{Does compiling change representations?}
% Summary 
The learning dynamics suggest \t{Model F} is using the compiled attention head
to quickly learn the conditional reverse task. To check, we can observe whether
there is task-relevant structure in the representations proceeding
the compiled attention head versus the non-compiled attention heads. 
% Explain Figure
Fig. \ref{fig:pca} displays the representations of each model projected onto
their first three principal components, extracted using PCA. Row \t{R2} depicts
the residual stream after the attention layer, and \t{A} the outputs of the non-compiled
attention heads. For clarity, only representations from position 16 are shown, and on a fragment
of the alphabet. The points are labeled by target token identity, and use an \t{X} marker when
the first token in the input sequence is a vowel. 
% What is interesting
Crucially, there appears to be no task-relevant
structure produced by non-compiled attention heads \t{A} for \t{Model F}.

\begin{figure}[t] 
  \centering 
  \includegraphics[width=.85\linewidth]{../experiments/cond_reverse/data/plots/pca_res_step5000_lr0.01_bs128_seed15.pdf}
  \vspace{.3em}
  \caption{PCA on Representations}
  \label{fig:pca}
  \vspace{-1em}
\end{figure}

\vspace{-.25em}
\subsection{Does compiling change the necessary representations?}

\begin{wrapfigure}{r}{0.4\linewidth}
  \vspace{-.5em}
  \centering
  \includegraphics[width=\linewidth]{../experiments/cond_reverse/data/plots/ablations_step5000_lr0.01_bs128.pdf}
  \caption{Resample Ablations}
  \vspace{-1em}
  \label{fig:ablations}
\end{wrapfigure}

To verify whether the compiled attention head is necessary for behavior, we resample ablate, for each relevant model,
the representations produced by the compiled attention head \t{R} and non-compiled attention heads \t{A}. Resampling
ablation replaces representations with representions generated by another input \citep{chan2022causal}. 
Results are shown in Fig. \ref{fig:ablations}, where accuracy is next-token prediction across all seeds. The ablations
show \t{Model F} behaves essentially the same when ablating non-compiled attention heads, but other models rely
on them to varying degrees. Interestingly, \t{Model P}, is fairly robust to ablating its compiled attention head,
which contains a reversal algorithm. This shows that a transformer does not necessarily use any compiled program, even if
intuitively relevant to the task.

\subsection{Does compiling change the controllable representations?}

\begin{wrapfigure}[13]{r}{0.4\linewidth}
  \vspace{-1em}
  \centering
    \includegraphics[width=\linewidth]{../experiments/cond_reverse/data/plots/steering_modelf_step5000_lr0.01_bs128.pdf}
  \caption{Steering Privileged Bases}
  \label{fig:steering}
\end{wrapfigure}

An interesting feature of \t{Model F} is that we have a rich understanding of the compiled attention head,
including what coordinates of which representations are \textit{privileged}, whether they denote a human-intelligible
concept \citep{elhage2023privileged}. As a final test of whether the transformer is using \t{Model F} as intended, we test whether 
gradient descent has preserved these privileged bases. The $\cj$\! program compiled into \t{Model F} represents whether 
a character is a vowel in its first coordinate, and all other coordinates represent non-vowelness. By scaling these coordinates we should be able
to reliably induce reversal (\t{Only Rev}) or identity (\t{Only Id})---no matter if the first character is a vowel \t{V} or a consonant \t{C}. Fig. \ref{fig:steering}
suggests these privileged bases are preserved.
